\chapter{Heavy Periods (Menorrhagia)}

\section*{Definition}
Excessive or prolonged menstrual bleeding defined as soaking through one or more pads/tampons every hour for several consecutive hours, bleeding lasting more than 7 days, or passing large blood clots.

\section*{What Happens in Your Body}
Menorrhagia results from hormonal imbalance (estrogen-progesterone ratio disruption), uterine structural abnormalities (fibroids, polyps, adenomyosis), coagulation disorders, or endometrial dysfunction. Anovulatory cycles common in perimenopause and adolescence lead to unopposed estrogen stimulation causing thickened endometrium and heavy bleeding.

\section*{Causes}
\begin{itemize}
  \item Uterine fibroids (leiomyomas)
  \item Hormonal imbalance (anovulation, perimenopause)
  \item Adenomyosis
  \item Endometrial polyps
  \item Coagulation disorders (von Willebrand disease, platelet disorders)
  \item Pelvic inflammatory disease
  \item Endometriosis
  \item Thyroid dysfunction (hypothyroidism)
  \item Intrauterine device (IUD) side effects
  \item Endometrial hyperplasia or malignancy
\end{itemize}

\section*{When to See a Doctor}
See your doctor if:
\begin{itemize}
  \item Symptoms are severe or getting worse over time
  \item Profuse bleeding soaking more than 1 pad/tampon per hour, bleeding with dizziness or fatigue, or between-period bleeding
  \item You have other concerning symptoms such as fever, weight loss, or bleeding
\end{itemize}

\section*{Related Symptoms}
\begin{itemize}
  \item Fatigue (anemia)
  \item Abdominal cramps
  \item Irregular periods
  \item Back pain
  \item Pale skin
\end{itemize}
\textbf{Important:} If this symptom persists, your doctor will check for serious underlying causes.

\section*{Self-Care / Home Management}
\begin{itemize}
  \item Rest and avoid activities that make it worse.
  \item Maintain adequate hydration and a balanced diet.
  \item Over-the-counter remedies may provide symptomatic relief (consult a pharmacist or doctor).
  \item Monitor symptoms and seek professional advice if they persist or worsen.
  \item Avoid self-medicating with prescription medications without medical guidance.
\end{itemize}

\section*{How Your Doctor Will Evaluate You}
\begin{itemize}
  \item A thorough history and physical examination are the first steps in evaluation.
  \item Laboratory tests (blood work, urinalysis) may be ordered based on clinical suspicion.
  \item Imaging studies (X-ray, ultrasound, CT, MRI) may be indicated when structural pathology is suspected.
  \item Specialist referral is arranged when the diagnosis remains uncertain or requires specific expertise.
\end{itemize}

\section*{Treatment Options}
\begin{itemize}
  \item Treatment targets the underlying cause identified through diagnostic evaluation.
  \item Supportive care and symptom management are provided while awaiting definitive therapy.
  \item Medications, lifestyle modifications, and physical therapies may be prescribed as appropriate.
  \item Follow-up ensures treatment effectiveness and allows timely adjustments to the care plan.
\end{itemize}

\clearpage